\documentclass[UTF8,12pt]{ctexart}
\usepackage[a4paper,margin=24mm]{geometry}
\usepackage{amsmath,amssymb,bm,graphicx,adjustbox}
\newcommand{\fitmath}[1]{\begin{adjustbox}{max width=\linewidth}$\displaystyle #1$\end{adjustbox}}
\setlength{\parindent}{0pt}
\setlength{\parskip}{.7em}
\sloppy
\allowdisplaybreaks
\begin{document}
\section*{1989 年考研数学三 · 真题}
\subsection*{1989年全国硕士研究生招生考试试题}

\subsection*{（试卷IV）}

\subsection*{一、填空题（本题共5小题，每小题3分，满分15分）}

（1）曲线 \(\displaystyle y=x+\sin^2x\) 在点 \(\displaystyle (\dfrac{\displaystyle \pi}{\displaystyle 2},\allowbreak 1+\dfrac{\displaystyle \pi}{\displaystyle 2})\) 处的切线方程是\_\_\_\_\_\_。


（2）幂级数 \(\displaystyle \sum_{n=0}^{\infty}\dfrac{\displaystyle x^n}{\displaystyle \sqrt{n+1}}\) 的收敛域是\_\_\_\_\_\_。


（3）若齐次线性方程组


\[\fitmath{\displaystyle \begin{cases}\lambda x_1+x_2+x_3=0,\\x_1+\lambda x_2+x_3=0,\\x_1+x_2+x_3=0\end{cases}}\]


只有零解，则 \(\displaystyle \lambda\) 应满足的条件是\_\_\_\_\_\_。


（4）设随机变量 \(\displaystyle X\) 的分布函数为


\[\fitmath{\displaystyle F(x)=\begin{cases}0,&x<0,\\A\sin x,&0\le x\le\dfrac{\displaystyle \pi}{\displaystyle 2},\\1,&x>\dfrac{\displaystyle \pi}{\displaystyle 2},\end{cases}}\]


则 \(\displaystyle A=\)\_\_\_\_\_\_，\(\displaystyle P\{|X|<\dfrac{\displaystyle \pi}{\displaystyle 6}\}=\)\_\_\_\_\_\_。


（5）设随机变量 \(\displaystyle X\) 的数学期望 \(\displaystyle E(X)=\mu\)，方差 \(\displaystyle D(X)=\sigma^2\)，则由切比雪夫（Chebyshev）不等式，有 \(\displaystyle P\{|X-\mu|\ge3\sigma\}\le\)\_\_\_\_\_\_。


\subsection*{二、选择题（本题共5小题，每小题3分，满分15分）}

（1）设 \(\displaystyle f(x)=2^x+3^x-2\)，则当 \(\displaystyle x\to0\) 时（　　）


（A）\(\displaystyle f(x)\) 是 \(\displaystyle x\) 的等价无穷小量。 （B）\(\displaystyle f(x)\) 与 \(\displaystyle x\) 是同阶但非等价无穷小量。 （C）\(\displaystyle f(x)\) 是比 \(\displaystyle x\) 较高阶的无穷小量。 （D）\(\displaystyle f(x)\) 是比 \(\displaystyle x\) 较低阶的无穷小量。


（2）在下列等式中，正确的结果是（　　）


（A）\(\displaystyle \int f\prime(x)\,dx=f(x)\)。 （B）\(\displaystyle \int df(x)=f(x)\)。 （C）\(\displaystyle \dfrac{\displaystyle d}{\displaystyle dx}\displaystyle \int f(x)\,dx=f(x)\)。 （D）\(\displaystyle d\displaystyle \int f(x)\,dx=f(x)\)。


（3）设 \(\displaystyle A\) 为 \(\displaystyle n\) 阶方阵且 \(\displaystyle |A|=0\)，则（　　）


（A）\(\displaystyle A\) 中必有两行（列）的元素对应成比例。 （B）\(\displaystyle A\) 中任意一行（列）向量是其余各行（列）向量的线性组合。 （C）\(\displaystyle A\) 中必有一行（列）向量是其余各行（列）向量的线性组合。 （D）\(\displaystyle A\) 中至少有一行（列）的元素全为0。


（4）设 \(\displaystyle A\) 和 \(\displaystyle B\) 都是 \(\displaystyle n\times n\) 矩阵，则必有（　　）


（A）\(\displaystyle |A+B|=|A|+|B|\)。 （B）\(\displaystyle AB=BA\)。 （C）\(\displaystyle |AB|=|BA|\)。 （D）\(\displaystyle (A+B)^{-1}=A^{-1}+B^{-1}\)。


（5）以 \(\displaystyle A\) 表示事件“甲种产品畅销，乙种产品滞销”，则其对立事件 \(\displaystyle \bar A\) 为（　　）


（A）“甲种产品滞销，乙种产品畅销”。 （B）“甲、乙两种产品均畅销”。 （C）“甲种产品滞销”。 （D）“甲种产品滞销或乙种产品畅销”。


\subsection*{三、计算题（本题共3小题，每小题5分，满分15分）}

（1）求极限 \(\displaystyle \lim_{x\to\infty}(\sin\dfrac{\displaystyle 1}{\displaystyle x}+\cos\dfrac{\displaystyle 1}{\displaystyle x})^x\)。


（2）已知 \(\displaystyle z=f(u,\allowbreak v),\allowbreak \ u=x+y,\allowbreak \ v=xy\)，且 \(\displaystyle f(u,\allowbreak v)\) 的二阶偏导数都连续，求 \(\displaystyle \dfrac{\displaystyle \partial^2z}{\displaystyle \partial x\partial y}\)。


（3）求微分方程 \(\displaystyle y\prime\prime+5y\prime+6y=2e^{-x}\) 的通解。


\subsection*{四、（本题满分9分）}

设某厂家打算生产一批商品投放市场。已知该商品的需求函数为 \(\displaystyle p=p(x)=10e^{-x/2}\)，且最大需求量为6，其中 \(\displaystyle x\) 表示需求量，\(\displaystyle p\) 表示价格。


（1）求该商品的收益函数和边际收益函数；（2）求使收益最大的产量、最大收益和相应的价格；（3）画出收益函数的图形。


\subsection*{五、（本题满分9分）}

已知函数


\[\fitmath{\displaystyle f(x)=\begin{cases}x,&0\le x\le1,\\2-x,&1<x\le2.\end{cases}}\]


计算下列各题：


（1）\(\displaystyle S_0=\displaystyle \int_0^2f(x)e^{-x}\,dx\)；（2）\(\displaystyle S_1=\displaystyle \int_2^4f(x-2)e^{-x}\,dx\)；（3）\(\displaystyle S_n=\displaystyle \int_{2n}^{2n+2}f(x-2n)e^{-x}\,dx\;(n=2,\allowbreak 3,\allowbreak \ldots)\)；（4）\(\displaystyle S=\displaystyle \sum_{n=0}^{\infty}S_n\)。


\subsection*{六、（本题满分6分）}

假设函数 \(\displaystyle f(x)\) 在 \(\displaystyle [a,\allowbreak b]\) 上连续，在 \(\displaystyle (a,\allowbreak b)\) 内可导，且 \(\displaystyle f\prime(x)\le0\)。记 \(\displaystyle F(x)=\dfrac{\displaystyle 1}{\displaystyle x-a}\displaystyle \int_a^x f(t)\,dt\)，证明在 \(\displaystyle (a,\allowbreak b)\) 内，\(\displaystyle F\prime(x)\le0\)。


\subsection*{七、（本题满分5分）}

已知 \(\displaystyle X=AX+B\)，其中


\[\fitmath{\displaystyle A=\begin{pmatrix}0&1&0\\-1&1&1\\-1&0&-1\end{pmatrix},\quad B=\begin{pmatrix}1&-1\\2&0\\5&-3\end{pmatrix}.}\]


求矩阵 \(\displaystyle X\)。


\subsection*{八、（本题满分6分）}

设 \(\displaystyle \boldsymbol\alpha_1=(1,\allowbreak 1,\allowbreak 1),\allowbreak \boldsymbol\alpha_2=(1,\allowbreak 2,\allowbreak 3),\allowbreak \boldsymbol\alpha_3=(1,\allowbreak 3,\allowbreak t)\)。问：（1）当 \(\displaystyle t\) 为何值时，向量组 \(\displaystyle \boldsymbol\alpha_1,\allowbreak \boldsymbol\alpha_2,\allowbreak \boldsymbol\alpha_3\) 线性无关？（2）当 \(\displaystyle t\) 为何值时，向量组线性相关？（3）当向量组线性相关时，将 \(\displaystyle \boldsymbol\alpha_3\) 表示为 \(\displaystyle \boldsymbol\alpha_1\) 和 \(\displaystyle \boldsymbol\alpha_2\) 的线性组合。


\subsection*{九、（本题满分5分）}

设


\[\fitmath{\displaystyle A=\begin{pmatrix}-1&2&2\\2&-1&-2\\2&-2&-1\end{pmatrix}.}\]


（1）试求矩阵 \(\displaystyle A\) 的特征值；（2）利用（1）的结果，求矩阵 \(\displaystyle E+A^{-1}\) 的特征值，其中 \(\displaystyle E\) 是3阶单位矩阵。


\subsection*{十、（本题满分7分）}

已知随机变量 \(\displaystyle X\) 和 \(\displaystyle Y\) 的联合概率密度为


\[\fitmath{\displaystyle f(x,y)=\begin{cases}e^{-(x+y)},&0<x<+\infty,\ 0<y<+\infty,\\0,&\text{其他},\end{cases}}\]


试求：（1）\(\displaystyle P\{X<Y\}\)；（2）\(\displaystyle E(XY)\)。


\subsection*{十一、（本题满分8分）}

设随机变量 \(\displaystyle X\) 在 \(\displaystyle [2,\allowbreak 5]\) 上服从均匀分布，现在对 \(\displaystyle X\) 进行三次独立观测，试求至少有两次观测值大于3的概率。


\subsection*{（试卷V）}

\subsection*{一、填空题（本题共5小题，每小题3分，满分15分）}

（1）【同试卷IV 第一、（1）题】


（2）某商品的需求量 \(\displaystyle Q\) 与价格 \(\displaystyle P\) 的函数关系为 \(\displaystyle Q=aP^b\)，其中 \(\displaystyle a\) 和 \(\displaystyle b\) 为常数，且 \(\displaystyle a\ne0\)，则需求量对价格 \(\displaystyle P\) 的弹性是\_\_\_\_\_\_。


（3）行列式


\[\fitmath{\displaystyle \begin{vmatrix}1&-1&1&x-1\\1&-1&x+1&-1\\1&x-1&1&-1\\x+1&-1&1&-1\end{vmatrix}=\underline{\qquad}.}\]


（4）设随机变量 \(\displaystyle X_1,\allowbreak X_2,\allowbreak X_3\) 相互独立，其中 \(\displaystyle X_1\) 在 \(\displaystyle [0,\allowbreak 6]\) 上服从均匀分布，\(\displaystyle X_2\) 服从正态分布 \(\displaystyle N(0,\allowbreak 2^2)\)，\(\displaystyle X_3\) 服从参数为 \(\displaystyle \lambda=3\) 的泊松分布。记 \(\displaystyle Y=X_1-2X_2+3X_3\)，则 \(\displaystyle D(Y)=\)\_\_\_\_\_\_。


（5）【同试卷IV 第一、（4）题】


\subsection*{二、选择题（本题共5小题，每小题3分，满分15分）}

（1）【同试卷IV 第二、（1）题】


（2）【同试卷IV 第二、（2）题】


（3）【同试卷IV 第二、（3）题】


（4）设 \(\displaystyle n\) 元齐次线性方程组 \(\displaystyle AX=0\) 的系数矩阵 \(\displaystyle A\) 的秩为 \(\displaystyle r\)，则 \(\displaystyle AX=0\) 有非零解的充分必要条件是（　　）


（A）\(\displaystyle r=n\)。 （B）\(\displaystyle r<n\)。 （C）\(\displaystyle r\ge n\)。 （D）\(\displaystyle r>n\)。


（5）【同试卷IV 第二、（5）题】


\subsection*{三、（本题共4小题，每小题5分，共20分）}

（1）求极限 \(\displaystyle \lim_{x\to+\infty}(x+e^x)^{1/x}\)。


（2）已知 \(\displaystyle z=a^{\sqrt{x^2-y^2}}\)，其中 \(\displaystyle a>0,\allowbreak a\ne1\)，求 \(\displaystyle dz\)。


（3）求不定积分 \(\displaystyle \int\dfrac{\displaystyle x+\ln(1-x)}{\displaystyle x^2}\,dx\)。


（4）求二重积分 \(\displaystyle \iint_D\dfrac{\displaystyle 1-x^2-y^2}{\displaystyle 1+x^2+y^2}\,dx\,dy\)，其中 \(\displaystyle D\) 是 \(\displaystyle x^2+y^2=1,\allowbreak x=0\) 和 \(\displaystyle y=0\) 所围成的区域在第I象限的部分。


\subsection*{四、（本题满分6分）}

已知某企业的总收入函数为 \(\displaystyle R=26x-2x^2-4x^3\)，总成本函数为 \(\displaystyle C=8x+x^2\)，其中 \(\displaystyle x\) 表示产品的产量，求利润函数，边际收入函数，边际成本函数，以及企业获得最大利润时的产量和最大利润。


\subsection*{五、（本题满分12分）}

已知函数 \(\displaystyle y=\dfrac{\displaystyle 2x^2}{\displaystyle (1-x)^2}\)，试求其单调区间、极值点及图形的凹凸性、拐点和渐近线，并画出函数的图形。


\subsection*{六、（本题满分5分）【同试卷IV 第七题】}

\subsection*{七、（本题满分6分）}

讨论向量组 \(\displaystyle \boldsymbol\alpha_1=(1,\allowbreak 1,\allowbreak 0),\allowbreak \boldsymbol\alpha_2=(1,\allowbreak 3,\allowbreak -1),\allowbreak \boldsymbol\alpha_3=(5,\allowbreak 3,\allowbreak t)\) 的线性相关性。


\subsection*{八、（本题满分5分）【同试卷IV 第九题】}

\subsection*{九、（本题满分8分）}

已知随机变量 \(\displaystyle X\) 和 \(\displaystyle Y\) 的联合概率分布为


\[\fitmath{\displaystyle \begin{array}{c|cccccc}(x,y)&(0,0)&(0,1)&(1,0)&(1,1)&(2,0)&(2,1)\\\hline P\{X=x,Y=y\}&0.10&0.15&0.25&0.20&0.15&0.15\end{array}}\]


试求：（1）\(\displaystyle X\) 的概率分布；（2）\(\displaystyle X+Y\) 的概率分布；（3）\(\displaystyle Z=\sin\dfrac{\displaystyle \pi(X+Y)}{\displaystyle 2}\) 的数学期望。


\subsection*{十、（本题满分8分）}

某仪器装有三只独立工作的同型号电子元件，其寿命（单位：小时）都服从同一指数分布，概率密度为


\[\fitmath{\displaystyle f(x)=\begin{cases}\dfrac{\displaystyle 1}{\displaystyle 600}e^{-x/600},&x>0,\\0,&x\le0,\end{cases}}\]


试求：在仪器使用的最初200小时内，至少有一只电子元件损坏的概率 \(\displaystyle \alpha\)。

\end{document}
